\begin{textsample}{NEG Dim 6 – global\_south\_2024\_03 – Score 1.00 – global\_south\_2024\_03/106330935.txt}  \label{ex:f6_neg_412}
Palestinians attend Friday prayers near the ruins of a mosque destroyed in Israeli strikes, amid the ongoing conflict between Israel and Hamas, in Rafah in the southern Gaza Strip Reuters After the strike on the Palestinians waiting for aid in Gaza, many countries have raised questions against Israel seeking an explanation behind the attack.
France, Germany, Lebanon and Brazil were among the countries that sought an investigation into the attacks.
Though Israel has denied the strikes on crowds waiting for aid, the death of 112 Gazans after the strikes is clearly not something that can be ignored.
In a post on social media, Germany ' s foreign minister Annalena Baerbock has demanded that Israel ' s army must"fully explain"how Palestinians were killed while they were gathered to receive humanitarian aid."People wanted relief supplies for themselves and their families and found themselves dead.
The reports from Gaza shock me.
The Israeli army must fully explain how the mass panic and shooting could have happened.
My condolences go out to the families of the victims,"wrote @ @ @ @ @ @ @ @ @ @ than to living.
More humanitarian aid needs to come in, she added.
President of the European Commission Ursula von der Leyen said she is deeply disturbed by the images from Gaza."Every effort must be made to investigate what happened and ensure transparency Humanitarian aid is a lifeline for those in need and access to it must be ensured,"she added.
Similarly, France ' s Foreign Minister St? phane S? journ? also called for an independent investigation into the deaths of Palestinians."We will ask for explanations, and there will have to be an independent probe to determine what happened.
France calls things by their name.
This applies when we designate Hamas as a terrorist group, but we must also call things by their name when there are atrocities in Gaza,"said S? journ?.
Lebanon ' s foreign ministry also echoed the call for an investigation.
While condemning the attack, the ministry said the attacks were"deliberate"targeting defenseless Palestinians. @ @ @ @ @ @ @ @ @ @ civilians."Such loss of civilian lives and the larger humanitarian situation in Gaza continues to be a cause for extreme concern,"said foreign ministry in a statement.
Brazil called Israel ' s action as"unethical and illegal"."Humanity is failing the civilians of Gaza.
And it ' s time to prevent further massacres,"Brazil ' s foreign ministry said in a statement.
It also added that it is up to the international community to call a halt to it.
Meanwhile, at least 30,228 Palestinians were killed and 71,377 have been injured in Israel ' s military offensive in Gaza.

% matched lemmas: 
\end{textsample}
